\chapter{Complexity — Logspace Uniform Canon}
%%% generated by the blueprint pipeline from TCSlib.Complexity.CircuitComplexity.LogspaceUniformCanon
\section{Overview}

Declarations of the Lean module \texttt{TCSlib.Complexity.CircuitComplexity.LogspaceUniformCanon}.

\section{Declarations}

\begin{definition}[Canonical gate arguments]
\lean{BoolCircuit.DAGGate.canon}
\label{BoolCircuit.DAGGate.canon}
\leanok
\uses{BoolCircuit.DAGGate}
A gate’s canonical form retains its gate type and replaces its argument list by the same
arguments in increasing order, with repetitions removed.
\end{definition}

\begin{lemma}[Bound on an element of a list]
\lean{BoolCircuit.DAGGate.le_foldr_max}
\label{BoolCircuit.DAGGate.le_foldr_max}
\leanok
\difficulty{4}
If a natural number $a$ belongs to a list $l$ of natural numbers, then $a$ is at most
the result of folding the maximum operation over $l$, starting from $0$.
\end{lemma}

\begin{lemma}[Argument membership under canonicalization]
\lean{BoolCircuit.DAGGate.mem_canon_args}
\label{BoolCircuit.DAGGate.mem_canon_args}
\leanok
\uses{BoolCircuit.DAGGate, BoolCircuit.DAGGate.canon, BoolCircuit.DAGGate.le_foldr_max}
\difficulty{3}
Let $g$ be a directed acyclic graph gate and let $a$ be a natural number. The number $a$
occurs in the argument list of the canonical form of $g$ if and only if it occurs in the
argument list of $g$.
\end{lemma}

\begin{lemma}[Canonical form preserves gate kind]
\lean{BoolCircuit.DAGGate.canon_kind}
\label{BoolCircuit.DAGGate.canon_kind}
\leanok
\uses{BoolCircuit.DAGGate, BoolCircuit.DAGGate.canon}
\difficulty{1}
For every Boolean gate, its canonical form has the same gate kind as the original gate.
\end{lemma}

\begin{lemma}[Canonical arguments are sorted]
\lean{BoolCircuit.DAGGate.canon_sorted}
\label{BoolCircuit.DAGGate.canon_sorted}
\leanok
\uses{BoolCircuit.DAGGate, BoolCircuit.DAGGate.canon}
\difficulty{1}
For every Boolean gate, the argument indices of its canonical form are strictly
increasing.
\end{lemma}

\begin{lemma}[Evaluation preserved by canonicalization]
\lean{BoolCircuit.DAGGate.eval_canon}
\label{BoolCircuit.DAGGate.eval_canon}
\leanok
\uses{BoolCircuit.DAGGate, BoolCircuit.DAGGate.canon, BoolCircuit.DAGGate.canon_kind, BoolCircuit.DAGGate.eval}
\difficulty{4}
For every gate $g$ and every finite list $v$ of Boolean values, the canonical form of
$g$ has the same evaluation on $v$ as $g$ does.
\end{lemma}

\begin{lemma}[Invariance under gate relabeling]
\lean{BoolCircuit.runWith_map}
\label{BoolCircuit.runWith_map}
\leanok
\uses{BoolCircuit.DAGGate, BoolCircuit.runWith, BoolCircuit.runWith_cons}
\difficulty{4}
Let $\beta$ be a type, let $f$ assign an output in $\beta$ to each gate and list of
previous outputs, and let $h$ map gates to gates. Suppose $f(h(g),v)=f(g,v)$ for every
gate $g$ and output list $v$. Then running the gates obtained by applying $h$ to each
gate in a list $gs$, starting from an initial output list, gives the same output list as
running $gs$ from that initial list.
\end{lemma}

\begin{definition}[Canonical circuit]
\lean{BoolCircuit.DAGCircuit.canon}
\label{BoolCircuit.DAGCircuit.canon}
\leanok
\uses{BoolCircuit.DAGCircuit, BoolCircuit.DAGGate.canon, BoolCircuit.DAGGate.mem_canon_args}
For a circuit with $n$ inputs, its canonical form keeps the output vertex and replaces
each gate by its canonical form, so the gates remain in their original order. The
resulting circuit has the same input count and valid gate and output indices as the
original.
\end{definition}

\begin{lemma}[Size preserved by canonicalization]
\lean{BoolCircuit.DAGCircuit.size_canon}
\label{BoolCircuit.DAGCircuit.size_canon}
\leanok
\uses{BoolCircuit.DAGCircuit, BoolCircuit.DAGCircuit.canon, BoolCircuit.DAGCircuit.size}
\difficulty{1}
Let $C$ be a finite Boolean circuit with $n$ input vertices, gates in topological order,
and a designated output vertex. Its canonicalized circuit has the same size as $C$.
\end{lemma}

\begin{lemma}[Evaluation of the canonical circuit]
\lean{BoolCircuit.DAGCircuit.eval_canon}
\label{BoolCircuit.DAGCircuit.eval_canon}
\leanok
\uses{BoolCircuit.DAGCircuit, BoolCircuit.DAGCircuit.canon, BoolCircuit.DAGCircuit.eval, BoolCircuit.DAGCircuit.values, BoolCircuit.DAGGate.eval_canon, BoolCircuit.runWith_map}
\difficulty{3}
For any directed acyclic Boolean circuit $C$ with $n$ inputs, its canonical form has the
same evaluation as $C$ on every input assignment.
\end{lemma}

\begin{lemma}[Vertex types are preserved by canonicalization]
\lean{BoolCircuit.DAGCircuit.vtype_canon}
\label{BoolCircuit.DAGCircuit.vtype_canon}
\leanok
\uses{BoolCircuit.DAGCircuit, BoolCircuit.DAGCircuit.canon, BoolCircuit.DAGCircuit.vtype}
\difficulty{4}
For every Boolean directed acyclic circuit $C$ and every vertex index $i$, the gate type
at vertex $i$ in the canonicalized circuit equals the gate type at vertex $i$ in $C$.
\end{lemma}

\begin{lemma}[Canonicalization preserves edges]
\lean{BoolCircuit.DAGCircuit.edge_canon}
\label{BoolCircuit.DAGCircuit.edge_canon}
\leanok
\uses{BoolCircuit.DAGCircuit, BoolCircuit.DAGCircuit.Edge, BoolCircuit.DAGCircuit.canon, BoolCircuit.DAGGate.canon, BoolCircuit.DAGGate.mem_canon_args}
\difficulty{5}
Let $C$ be a directed acyclic Boolean circuit on $n$ input vertices, and let $i,j$ be
vertices. There is an edge from $i$ to $j$ in the canonical form of $C$ if and only if
there is an edge from $i$ to $j$ in $C$.
\end{lemma}

\begin{lemma}[Canonical form of a circuit]
\lean{BoolCircuit.DAGCircuit.isCanonical_canon}
\label{BoolCircuit.DAGCircuit.isCanonical_canon}
\leanok
\uses{BoolCircuit.DAGCircuit, BoolCircuit.DAGCircuit.IsCanonical, BoolCircuit.DAGCircuit.canon, BoolCircuit.DAGCircuit.size, BoolCircuit.DAGCircuit.size_canon, BoolCircuit.DAGGate.canon_sorted}
\difficulty{4}
Let $C$ be a Boolean circuit on $n$ inputs whose output vertex is the last vertex, so
that its output index plus one equals its size. Then the circuit obtained by
canonicalizing $C$ has, at each gate, its input indices in strictly increasing order,
and its output vertex is the last vertex.
\end{lemma}

\begin{definition}[Canonical circuit family]
\lean{BoolCircuit.DAGCircuitFamily.canon}
\label{BoolCircuit.DAGCircuitFamily.canon}
\leanok
\uses{BoolCircuit.DAGCircuit.canon, BoolCircuit.DAGCircuitFamily}
A circuit family assigns to each input length $n$ the canonical form of its circuit for
inputs of length $n$.
\end{definition}

\begin{lemma}[Canonical circuits preserve logspace adjacency]
\lean{BoolCircuit.DAGCircuitFamily.hasLogspaceAdjacency_canon}
\label{BoolCircuit.DAGCircuitFamily.hasLogspaceAdjacency_canon}
\leanok
\uses{BoolCircuit.DAGCircuitFamily, BoolCircuit.DAGCircuitFamily.HasLogspaceAdjacency, BoolCircuit.DAGCircuitFamily.canon, BoolCircuit.DAGCircuitFamily.edgeLang, BoolCircuit.DAGCircuitFamily.sizeLang, BoolCircuit.DAGCircuitFamily.typeLang}
\difficulty{5}
Let $C$ be a family of Boolean circuits, and suppose its size language, each of its
gate-type languages, and its edge language are decidable in logarithmic space, and the
family has polynomial size. Then the family obtained by replacing each circuit with its
canonical form also has polynomial size and has size, gate-type, and edge languages
decidable in logarithmic space.
\end{lemma}

\begin{theorem}[Canonical uniform circuit family]
\lean{BoolCircuit.DAGCircuitFamily.HasLogspaceAdjacency.exists_canonical_isLogspaceUniform}
\label{BoolCircuit.DAGCircuitFamily.HasLogspaceAdjacency.exists_canonical_isLogspaceUniform}
\leanok
\uses{BoolCircuit.DAGCircuit.eval, BoolCircuit.DAGCircuit.eval_canon, BoolCircuit.DAGCircuit.isCanonical_canon, BoolCircuit.DAGCircuit.size, BoolCircuit.DAGCircuitFamily, BoolCircuit.DAGCircuitFamily.HasLogspaceAdjacency, BoolCircuit.DAGCircuitFamily.IsCanonical, BoolCircuit.DAGCircuitFamily.IsLogspaceUniform, BoolCircuit.DAGCircuitFamily.canon, BoolCircuit.DAGCircuitFamily.hasLogspaceAdjacency_canon, BoolCircuit.DAGCircuitFamily.isLogspaceUniform_iff}
\difficulty{4}
Let $C$ be a family of directed acyclic Boolean circuits with logspace-decidable
adjacency information, and suppose that for every input length $n$ the output vertex of
$C_n$ satisfies $\mathrm{output}(C_n)+1=\mathrm{size}(C_n)$. Then there is a canonical
circuit family $C'$ such that $C'_n$ and $C_n$ compute the same Boolean function for
every $n$, and $C'$ is logspace-uniform.
\end{theorem}
